% FILE: 02_lineage_and_context.tex
% Project: audits
% Role: completeness-audit-and-compliance-verification-corpus

\section{Lineage and Context}
\label{sec:audits-lineage}

\subsection{2016 Language-Model Lesson}
\label{sec:audits-lineage-lm}

The 2016 language-model experiment established the founding negative result of
this research program: local text imitation does not conserve consequence. A
system can produce outputs that are locally indistinguishable from correct
outputs while failing to preserve the causal and temporal structure that makes
those outputs meaningful as process evidence.

This lesson directly motivates the audits corpus. A test suite that checks
whether \texttt{to\_json\_string()} returns a valid JSON string is performing
local text imitation~--- it verifies the form of the output but not its
consequence within the process-law chain. The \textsc{dto\_001} violation is not
a behavioral defect: the method returns a correct string. It is a
\emph{consequence defect}: the method name signals that structured
\texttt{Evidence<T>} has been collapsed to a casual carrier, destroying the
downstream receipting capability. Only a process-level audit~--- one that checks
the type-law surface, not the output surface~--- can detect this class of
defect.

The audits corpus is therefore the direct descendant of the 2016 lesson applied
to quality enforcement: conformance checking must operate at the level of
process law, not behavioral output.

\subsection{Enterprise Process Gap}
\label{sec:audits-lineage-enterprise}

The enterprise process gap that motivates the full research program is the
systematic inability of enterprise software systems to provide evidence that the
process they declared as their operating model is the process that actually
executed. Event logs exist, but they are not systematically derived from OTel
traces into object-centric OCEL format. Process models are declared in BPMN or
Petri net notation, but they are not routinely compared against discovered models
from the event evidence. When auditors ask whether a compliance process executed
correctly, the honest answer in most enterprise systems is: the code ran, the
test passed, but the process is unwitnessed.

The audits corpus addresses this gap within the research foundry itself. By
applying Van der Aalst conformance auditing to the foundry's own codebase~---
using deterministic shell scripts as the audit instrument and committed text
files as the receipt chain~--- the corpus demonstrates that the methodology
closes the enterprise process gap at the system level. The \texttt{ALIVE\_001}
adversarial audit certifying 567 commits across 12 corpus directories is the
receipt that the foundry's own process is audited and witnessed.

\subsection{Relationship to the Chatman Equation}
\label{sec:audits-lineage-chatman-equation}

The Chatman Equation, \(A = \mu(O)\), states that autonomic alignment \(A\) is
the measure of operator coverage \(\mu\) over the object space \(O\). In the
context of the audits corpus, this equation manifests as the requirement that
every audit result~--- every \textsc{pass} or \textsc{fail}~--- must be a
measure over a defined object space.

For the five deterministic shell audits, the object space is the
\texttt{wasm4pm-compat} crate surface: its method names, its feature
dependencies, its exported public items, its \texttt{Receiptable} trait
implementations. Each audit defines a measurement function \(\mu\) over this
surface and returns a scalar verdict. The graduation checkpoint
\texttt{PROCESS\_INTELLIGENCE\_COMPAT\_CONFORMANCE\_001} is admissible only when
\(\mu\) returns \textsc{pass} for all five measurement functions simultaneously.

The neuro-symbolic verification audit (\texttt{neuro-symbolic-verification.md})
references the Chatman Equation directly in the context of type-law coverage:
\emph{``By enforcing the Chatman Equation (\(A = \mu(O^*)\)), the system
guarantees that no adversarial process topology can escape explicit symbolic
resolution.''} This is an \textsc{author thesis} interpretation of the equation
applied to the audit domain~--- [AUTHOR THESIS].

\subsection{Relationship to Receipts and Replay}
\label{sec:audits-lineage-receipts}

The receipts-and-replay doctrine holds that no claim is admissible unless it is
accompanied by a cryptographic receipt that can be replayed to regenerate the
claim from first principles. The audits corpus operationalizes this doctrine at
two levels.

At the \emph{shell level}, each of the five deterministic audits is
re-executable: running \texttt{bash run\_audit\_1.sh} must produce byte-identical
output to the committed \texttt{audit-no-dto-flattening.txt} receipt. This
reproducibility constraint is enforced by the audit's determinism~--- all pattern
searches and structural traversals are deterministic Bash operations with no
random or time-dependent components.

At the \emph{cryptographic level}, the agent audit receipts are signed with
SHA-256 hashes and, for the \texttt{ALIVE\_001} certification, with BLAKE3 and
Ed25519. The genesis block SHA-256 hash
\texttt{787985dd49ff98f9803851093406bdfc2eb7ab5f71b374d01d5b9a0ea952fc26} and the
\texttt{ALIVE\_001} signature \texttt{BLAKE3(567\_commits\_15\_dirs) +
Ed25519(teamwork\_preview\_worker\_m5\_1)} constitute a cryptographic ledger that
is itself audited by the ledger custodian agent, whose \texttt{test\_sha256\_
ledger\_compliance} Rust integration test verifies cross-language equivalence of
the SHA-256 block chain between \texttt{experiments/visualizer/ledger.js} and
\texttt{sources/wasm4pm/src/crypto.rs}.
